\unitopen{31}
\begin{candos}
\item użyć najczęstszych czasowników ze stałym przyimkiem i właściwym przypadkiem;
\item zadać pytanie z \de{wo(r)-} (\de{Worauf wartest du?}) i odpowiedzieć z \de{da(r)-} (\de{Ich warte darauf});
\item odróżnić pytanie o osobę (\de{Auf wen?}) od pytania o rzecz (\de{Worauf?});
\item wykorzystać czasowniki z przyimkami w opisie pracy i spotkania.
\end{candos}

\section{Słownictwo}
\vocabtable{31}

\section{Czasownik i przyimek tworzą parę}
Przyimek po czasowniku często jest \emph{inny} niż w polskim, a razem z nim dziedziczysz przypadek. Uczysz się całych połączeń:

\begin{center}\small
\begin{tabular}{@{}lll@{}}
\toprule
\textsf{\textbf{Połączenie}} & \textsf{\textbf{Przypadek}} & \textsf{\textbf{Przykład}}\\
\midrule
\de{warten auf} & Akk & \de{Ich warte auf den Bus.}\\
\de{denken an} & Akk & \de{Ich denke an den Termin.}\\
\de{sich interessieren für} & Akk & \de{Er interessiert sich für KI.}\\
\de{sprechen über} & Akk & \de{Wir sprechen über das Projekt.}\\
\de{teilnehmen an} & Dat & \de{Ich nehme am Meeting teil.}\\
\de{abhängen von} & Dat & \de{Das hängt vom Budget ab.}\\
\de{sich beschäftigen mit} & Dat & \de{Sie beschäftigt sich mit Architektur.}\\
\de{fragen nach} & Dat & \de{Er fragt nach dem Preis.}\\
\de{Angst haben vor} & Dat & \de{Ich habe Angst vor der Prüfung.}\\
\bottomrule
\end{tabular}
\end{center}

\section{\de{da(r)-} i \de{wo(r)-}}
Gdy w zdaniu zamiast rzeczy stoi zaimek, nie mówisz \de{auf es}, tylko \de{darauf}. Dla \emph{rzeczy} używasz \de{da(r)-} +
przyimek (\de{darauf, daran, darüber, damit, dafür}), dla \emph{osób} przyimek z zaimkiem (\de{auf ihn, an sie, mit ihm}). Pytanie o rzecz:
\de{wo(r)-} + przyimek (\de{Worauf wartest du? Woran denkst du? Womit arbeitest du?}); o osobę: \de{Auf wen wartest du?} \de{Mit wem sprichst du?}
\emph{r} pojawia się przed samogłoską: \de{darauf}, \de{woran}, ale \de{damit}, \de{womit}.

\begin{falle}
Polskie „czekam na” to \de{warten auf} (nie \de{an}), „myślę o” to \de{denken an}, „mówię o” to \de{sprechen über},
„pytam o” to \de{fragen nach}. Błąd: \de{Ich warte für den Bus}. Drugi: \de{auf es} zamiast \de{darauf}. Trzeci:
\de{worauf} dla osoby: o osobie \de{auf wen}.
\end{falle}

\begin{szwajcaria}
W rozmowie o pracy Szwajcarzy używają często \de{sich kümmern um} (\de{Ich kümmere mich darum}, „zajmę się tym”) i \de{Rücksprache halten mit}
(skonsultować). Odpowiedź \de{Ich schaue dazu} lub \de{Ich kümmere mich darum} jest normalnym sposobem potwierdzenia, że bierzesz sprawę.
\end{szwajcaria}

\begin{praca}
\de{Ich kümmere mich um den Fehler.} \de{Wir diskutieren über die Architektur.} \de{Das hängt von der Last ab.}
\de{Kannst du an dem Workshop teilnehmen?} \de{Worüber sprechen wir heute?} \de{Darüber habe ich noch nicht nachgedacht.}
\end{praca}

\section{Dialogi}
\textbf{Dialog 1 --- spotkanie projektowe}
\begin{dialog}
\turn{A}{Worüber diskutieren wir heute?}{O czym dziś dyskutujemy?}
\turn{B}{Über die Migration. Das hängt vom Budget ab. Daran müssen wir denken.}{O migracji. To zależy od budżetu. O tym musimy pamiętać.}
\turn{A}{Wer kümmert sich um die Planung?}{Kto zajmie się planowaniem?}
\turn{B}{Anna. Sie beschäftigt sich schon seit Wochen damit.}{Anna. Zajmuje się tym od tygodni.}
\end{dialog}

\section{Ćwiczenia}
\begin{cwiczenia}
\exitem{Uzupełnij przyimek: \de{Ich warte \dots{} den Bus.}}{\de{auf}}
\exitem{Uzupełnij przyimek: \de{Ich denke \dots{} den Termin.}}{\de{an}}
\exitem{Uzupełnij przyimek: \de{Wir sprechen \dots{} das Projekt.}}{\de{über}}
\exitem{Uzupełnij przyimek: \de{Das hängt \dots{} Budget ab.} (z rodzajnikiem w jednym słowie)}{\de{vom} (\de{von} + \de{dem})}
\exitem{Uzupełnij: \de{Ich nehme \dots{} Meeting teil.} (z rodzajnikiem w jednym słowie)}{\de{am}}
\exitem{Zadaj pytanie o rzecz: \de{Ich warte auf den Bus.}}{\de{Worauf wartest du?}}
\exitem{Zadaj pytanie o osobę: \de{Ich warte auf meinen Chef.}}{\de{Auf wen wartest du?}}
\exitem{Zastąp zaimkiem: \de{Ich denke an den Termin.} $\rightarrow$ \de{Ich denke \dots{}}.}{\de{Ich denke daran.}}
\exitem{Zastąp zaimkiem: \de{Ich spreche mit dem Chef.} $\rightarrow$ \de{Ich spreche \dots{}}.}{\de{Ich spreche mit ihm.}}
\exitem{Popraw: \de{Ich warte für den Bus.}}{\de{Ich warte auf den Bus.}}
\exitem{Przetłumacz: „Zajmę się tym.”}{\de{Ich kümmere mich darum.}}
\exitem{Przetłumacz: „O czym rozmawiacie?”}{\de{Worüber sprecht ihr?}}
\end{cwiczenia}

\section{Trening słownictwa (generowany)}
\drill{31}

\weekend{Wypisz po dziesięć czasowników z przyimkami z tej jednostki i ułóż po jednym zdaniu o swojej pracy. Do każdego zdania
dopisz pytanie z \de{wo(r)-} i odpowiedź z \de{da(r)-}. Nagraj 90 sekund.}
